\documentclass[UTF8,11pt]{ctexart}
\usepackage[a4paper,margin=24mm]{geometry}
\usepackage{amsmath,amssymb,amsthm,mathtools,mathrsfs}
\usepackage[all]{xy}
\usepackage{xurl,hyperref,enumitem}
\usepackage{tikz}
\usetikzlibrary{arrows.meta,cd,decorations.markings,calc,patterns}
\hypersetup{hidelinks,breaklinks=true}
\setlength{\emergencystretch}{3em}
\allowdisplaybreaks
\newtheorem*{theorem}{Theorem}
\newtheorem*{thm}{Theorem}
\newtheorem*{lemma}{Lemma}
\newtheorem*{proposition}{Proposition}
\newtheorem*{prop}{Proposition}
\newtheorem*{corollary}{Corollary}
\newtheorem*{cor}{Corollary}
\newtheorem*{claim}{Claim}
\theoremstyle{definition}
\newtheorem*{definition}{Definition}
\newtheorem*{defn}{Definition}
\newtheorem*{remark}{Remark}
\newtheorem*{example}{Example}
\newcommand{\R}{\mathbb R}
\newcommand{\C}{\mathbb C}
\newcommand{\Z}{\mathbb Z}
\newcommand{\Q}{\mathbb Q}
\newcommand{\N}{\mathbb N}
\newcommand{\F}{\mathbb F}
\newcommand{\D}{\mathbb D}
\newcommand{\bB}{\mathbb B}
\newcommand{\bC}{\mathbb C}
\newcommand{\bR}{\mathbb R}
\newcommand{\bZ}{\mathbb Z}
\newcommand{\bN}{\mathbb N}
\newcommand{\bQ}{\mathbb Q}
\newcommand{\bD}{\mathbb D}
\newcommand{\bF}{\mathbb F}
\newcommand{\CP}{\mathbb{CP}}
\newcommand{\RP}{\mathbb{RP}}
\newcommand{\sO}{\mathscr O}
\newcommand{\sI}{\mathscr I}
\newcommand{\sF}{\mathscr F}
\newcommand{\sG}{\mathscr G}
\newcommand{\sH}{\mathscr H}
\newcommand{\sR}{\mathscr R}
\newcommand{\sM}{\mathscr M}
\newcommand{\sV}{\mathscr V}
\newcommand{\sU}{\mathscr U}
\DeclareMathOperator{\Hom}{Hom}
\DeclareMathOperator{\Ext}{Ext}
\DeclareMathOperator{\Tor}{Tor}
\DeclareMathOperator{\Spec}{Spec}
\DeclareMathOperator{\Proj}{Proj}
\DeclareMathOperator{\supp}{supp}
\DeclareMathOperator{\im}{im}
\DeclareMathOperator{\id}{id}
\DeclareMathOperator{\Tr}{Tr}
\DeclareMathOperator{\tr}{tr}
\DeclareMathOperator{\rank}{rank}
\DeclareMathOperator{\ord}{ord}
\DeclareMathOperator{\dist}{dist}
\DeclareMathOperator{\End}{End}
\DeclareMathOperator{\Aut}{Aut}
\DeclareMathOperator{\Gal}{Gal}
\DeclareMathOperator{\vol}{vol}
\DeclareMathOperator{\Cl}{Cl}
\DeclareMathOperator{\coker}{coker}
\DeclareMathOperator{\codim}{codim}
\DeclareMathOperator{\divergence}{div}
\DeclareMathOperator{\Ric}{Ric}
\DeclareMathOperator{\grad}{grad}
\DeclareMathOperator{\Hess}{Hess}
\newcommand{\abs}[1]{\left\lvert#1\right\rvert}
\newcommand{\sabs}[1]{\lvert#1\rvert}
\newcommand{\babs}[1]{\bigl\lvert#1\bigr\rvert}
\newcommand{\Babs}[1]{\Bigl\lvert#1\Bigr\rvert}
\newcommand{\bbabs}[1]{\biggl\lvert#1\biggr\rvert}
\newcommand{\BBabs}[1]{\Biggl\lvert#1\Biggr\rvert}
\newcommand{\norm}[1]{\left\lVert#1\right\rVert}
\newcommand{\snorm}[1]{\lVert#1\rVert}
\newcommand{\bnorm}[1]{\bigl\lVert#1\bigr\rVert}
\newcommand{\Bnorm}[1]{\Bigl\lVert#1\Bigr\rVert}
\newcommand{\bbnorm}[1]{\biggl\lVert#1\biggr\rVert}
\newcommand{\BBnorm}[1]{\Biggl\lVert#1\Biggr\rVert}
\newcommand{\widebar}[1]{\overline{#1}}
\newcommand{\nicefrac}[2]{\frac{#1}{#2}}
\newcommand{\linnprod}[2]{\langle#1,#2\rangle}
\newcommand{\myindex}[1]{#1}
\newcommand{\glsadd}[1]{}
\newcommand{\avoidbreak}{}
\newcommand{\sourceref}[1]{\textnormal{[原文：\nolinkurl{#1}]}}
\newcommand{\thmref}[1]{\sourceref{#1}}
\newcommand{\lemmaref}[1]{\sourceref{#1}}
\newcommand{\propref}[1]{\sourceref{#1}}
\newcommand{\corref}[1]{\sourceref{#1}}
\newcommand{\defnref}[1]{\sourceref{#1}}
\newcommand{\exampleref}[1]{\sourceref{#1}}
\newcommand{\exerciseref}[1]{\sourceref{#1}}
\newcommand{\figureref}[1]{\sourceref{#1}}
\newcommand{\sectionref}[1]{\sourceref{#1}}
\newcommand{\appendixref}[1]{\sourceref{#1}}
\newcommand{\stacksref}[2]{\href{https://stacks.math.columbia.edu/tag/#1}{#2 [#1]}}

\begin{document}
\section*{033\quad 有界圆形域双全纯映射的线性刚性}
\subsection*{来源与计数目标}
Ji\v{r}\'i Lebl, \emph{Tasty Bits of Several Complex Variables}。从作者公开的原生 TeX 正文摘录；获取日期：2026-09-28。使用实际访问的在线正文，不把工作分支冒称冻结的印刷版。
\par\url{https://raw.githubusercontent.com/jirilebl/scv/master/scv.tex}
\par\url{https://www.jirka.org/scv/}
原文位置：\emph{Cartan\textquotesingle s uniqueness theorem} 节的 Cartan 定理及圆形域推论，两份作者证明。
\par\textbf{本文件只计一个问题：}固定原点的有界圆形域双全纯映射必为线性；Cartan 唯一性作为核心支撑合并收录。
\subsection*{作者命题与人类证明}

\paragraph{Cartan's uniqueness theorem (support for the rigidity conclusion).}
\begin{thm}[Cartan]
Suppose $U\subset\C^n$ is a bounded domain, $a\in U$, $f:U\to U$ is a holomorphic mapping, $f(a)=a$, and $Df(a)$ is the identity. Then $f(z)=z$ for all $z\in U$.
\end{thm}
A polynomial $P:\C^n\to\C$ is homogeneous of degree $d$ if $P(sz)=s^dP(z)$ for all $s\in\C$ and $z\in\C^n$. A polynomial vector-valued mapping is homogeneous of degree $d$ if each component is. The power series at $a$ of a holomorphic $f$ is then written as $\sum_{m=0}^\infty f_m(z-a)$, where $f_m$ is a homogeneous polynomial of degree $m$.
\begin{proof}
Without loss of generality, assume $a=0$. Write $f$ as a power series at the origin in terms of homogeneous parts:
\[
f(z)=z+f_k(z)+\sum_{m=k+1}^\infty f_m(z)=z+f_k(z)+\text{higher-order terms},
\]
where $k\geq2$ is an integer such that $f_2(z)=f_3(z)=\cdots=f_{k-1}(z)=0$. The degree-one homogeneous part is simply the vector $z$, because the derivative of $f$ at the origin is the identity. Compose $f$ with itself $\ell$ times:
\[f^\ell(z)=\underbrace{f\circ f\circ\cdots\circ f}_{\ell\text{ times}}(z).\]
As $f(U)\subset U$, we have that $f^\ell$ is a holomorphic map of $U$ to $U$. As $U$ is bounded, there is an $M$ such that $\|z\|\leq M$ for all $z\in U$. Therefore, $\|f(z)\|\leq M$ for all $z\in U$, and $\|f^\ell(z)\|\leq M$ for all $z\in U$. Note that
\[f_k\bigl(f(z)\bigr)=f_k\bigl(z+\text{higher-order terms}\bigr)=f_k(z)+\text{higher-order terms}.\]
Therefore,
\begin{align*}
f^2(z)&=f\bigl(f(z)\bigr)=f(z)+f_k\bigl(f(z)\bigr)+\text{higher-order terms}\\
&=z+2f_k(z)+\text{higher-order terms}.
\end{align*}
Continuing this procedure,
\[f^\ell(z)=z+\ell f_k(z)+\text{higher-order terms}.\]
Suppose $\Delta_r(0)$ is a polydisc whose closure is in $U$. Via Cauchy estimates, for every multi-index $\alpha$ with $|\alpha|=k$,
\[
\frac{\alpha!}{r^\alpha}M\geq\left\|\frac{\partial^{|\alpha|}f^\ell}{\partial z^\alpha}(0)\right\|=\ell\left\|\frac{\partial^{|\alpha|}f}{\partial z^\alpha}(0)\right\|.
\]
The inequality holds for all $\ell\in\N$, and so $\frac{\partial^{|\alpha|}f}{\partial z^\alpha}(0)=0$. Therefore, $f_k\equiv0$. On the domain of convergence of the expansion, we get $f(z)=z$, as there is no other nonzero homogeneous part in the expansion of $f$. As $U$ is connected, the identity theorem says $f(z)=z$ for all $z\in U$.
\end{proof}

A circular domain is a domain $U\subset\C^n$ such that if $z\in U$, then $e^{i\theta}z\in U$ for all $\theta\in\R$.
\begin{cor}
Suppose $U,V\subset\C^n$ are bounded circular domains with $0\in U$, $0\in V$, and $f:U\to V$ is a biholomorphic map such that $f(0)=0$. Then $f$ is linear.
\end{cor}
\begin{proof}
The map $g(z)=f^{-1}\bigl(e^{-i\theta}f(e^{i\theta}z)\bigr)$ is an automorphism of $U$, and via the chain rule, $g'(0)=I$. Therefore, Cartan says that $f^{-1}\bigl(e^{-i\theta}f(e^{i\theta}z)\bigr)=z$, or in other words,
\[f(e^{i\theta}z)=e^{i\theta}f(z).\]
Write $f$ near zero as $f(z)=\sum_{m=1}^\infty f_m(z)$ where $f_m$ are homogeneous polynomials of degree $m$ (notice $f_0=0$). Then
\[
\sum_{m=1}^\infty e^{i\theta}f_m(z)=e^{i\theta}\sum_{m=1}^\infty f_m(z)=\sum_{m=1}^\infty f_m(e^{i\theta}z)=\sum_{m=1}^\infty e^{im\theta}f_m(z).
\]
By the uniqueness of the Taylor expansion, $e^{i\theta}f_m(z)=e^{im\theta}f_m(z)$, or $f_m(z)=e^{i(m-1)\theta}f_m(z)$, for all $m$, all $z$, and all $\theta$. If $m\neq1$, we obtain that $f_m\equiv0$, which proves the claim.
\end{proof}

\subsection*{整理说明（非作者正文）}
标准先修为多变量 Taylor 展开、逐分量 Cauchy 估计、复合映射链式法则和恒等定理。原书把向量值 Cauchy 估计列为习题并说明可逐分量使用普通估计，本条不将该估计伪称已经重证。齐次展开定义摘取作者原文；无关反例和自同构群计算练习不纳入计数。

难度在于选择迭代使最低非线性项的系数线性增长，再用解析有界性迫使其消失；继而构造与旋转的交换子，把解析唯一性转成齐次分量的权重限制。不是把“圆对称”直接当作线性结论。原定理和推论只合计一个问题。
\subsection*{可修改的简短标注（非作者正文）}
$c$：有界复域之间映射的线性刚性。

$a$：全纯迭代；Taylor 齐次分量；Cauchy 估计；旋转群作用；复合交换子；恒等定理。

\texttt{cross\_theory\_status = yes}。迭代的系数增长被解析估计限制，旋转作用再将唯一性转为不同齐次次数的权重分离，排除所有非线性分量。位置：Cartan 证明的迭代与 Cauchy 估计，以及推论中 $g$ 的构造。

全部标注均为非穷尽、可修改初稿；未标注不表示未使用。
\subsection*{许可与修改记录}
原数学正文作者：Ji\v{r}\'i Lebl。作者网站及源书声明允许 CC BY-SA 4.0 或 CC BY-NC-SA 4.0 双许可；本条选择 CC BY-SA 4.0，整理部分亦采用同一许可。\url{https://creativecommons.org/licenses/by-sa/4.0/}。2026-09-28：助手截取题证、调整独立排版，加入来源、范围与初稿标注；不暗示作者认可本次整理。保留的是所用 TeX 正文摘录，未将其称为整书原件。
\end{document}
